\section{Conclusiones}

\subsection{Conclusion general}
VotoSeguro ejemplifica cómo varios mecanismos criptográficos pueden colaborar en un protocolo de votación electrónica simulado. Argon2id protege las contraseñas almacenadas; RSA-OAEP y AES-GCM permiten transmitir un voto en forma cifrada; las transacciones SQLite impiden la duplicación por cuenta; SHA-256 detecta alteraciones del registro encadenado; y Ed25519 autentica el manifiesto final publicado.

La selección de algoritmos depende de su función: el cifrado simétrico es adecuado para el contenido de las boletas, el asimétrico sirve para transportar su clave, los hashes para registrar huellas y las firmas para autenticar resultados. Ningún método por sí solo resuelve todos los requisitos.

El prototipo también permite reconocer límites importantes: una autoridad honesta sigue siendo un supuesto de confianza, no se ha demostrado anonimato completo frente a correlación de metadatos, no existe un mecanismo matemático de escrutinio públicamente verificable y el sistema local no está preparado para una elección oficial. Por tanto, su propósito es \textbf{didáctico y experimental}, no electoral operativo.

\newpage
\subsection{Conclusiones individuales}
\textbf{Axel Ivan Verastegui Patino - 2091219}:
\newline
Al realizar este proyecto nos dimos cuenta de que crear un sistema de votación electrónica seguro es mucho más complicado de lo que parece. Al principio podría pensarse que basta con registrar usuarios y contar votos, pero existen muchos aspectos que debemos considerar, como evitar votos duplicados, proteger las contraseñas y verificar que los resultados no hayan sido alterados. 
Aunque logramos desarrollar un prototipo funcional, entendimos que un sistema utilizado en elecciones reales necesitaría muchas más pruebas, medidas de seguridad y supervisión antes de poder implementarse.\\

\textbf{Julio Cesar Puente Reynoso - 2115192}:
\newline
text.\\

\textbf{Cecilia Ximena Gonzales Chavez - 2069031}:
\newline
bruh.\\

\textbf{Diego Alejandro Puente Ruiz - 2095327}:
\newline
thfck.\\

\textbf{Sebastian Flores Gonzales - 2104564}:
\newline
 more xmore\\

\section*{Referencias bibliográficas}
\addcontentsline{toc}{section}{Referencias bibliográficas}
\begin{thebibliography}{99}
\bibitem{nistvvsg} NIST. \emph{Voluntary Voting System Guidelines: Principles and Guidelines}. \url{https://www.nist.gov/itl/voting/vvsg-principles-and-guidelines}
\bibitem{owasp} OWASP. \emph{Password Storage Cheat Sheet}. \url{https://cheatsheetseries.owasp.org/cheatsheets/Password_Storage_Cheat_Sheet.html}
\bibitem{mdn} MDN Web Docs. \emph{SubtleCrypto API}. \url{https://developer.mozilla.org/en-US/docs/Web/API/SubtleCrypto}
\bibitem{flask} Pallets. \emph{Flask Security Considerations}. \url{https://flask.palletsprojects.com/en/stable/web-security/}
\bibitem{rfc8017} Moriarty et al. \emph{RFC 8017: PKCS \#1 RSA Cryptography Specifications Version 2.2}. \url{https://www.rfc-editor.org/rfc/rfc8017}
\bibitem{nistgcm} NIST. \emph{SP 800-38D: Recommendation for Block Cipher Modes of Operation: GCM and GMAC}. \url{https://doi.org/10.6028/NIST.SP.800-38D}
\bibitem{rfc8032} Josefsson y Liusvaara. \emph{RFC 8032: Edwards-Curve Digital Signature Algorithm (EdDSA)}. \url{https://www.rfc-editor.org/rfc/rfc8032}
\end{thebibliography}
